\documentclass{article}

\usepackage[english]{babel}
\usepackage[letterpaper,top=2cm,bottom=2cm,left=3cm,right=3cm]{geometry}
\usepackage{amsmath}
\usepackage{graphicx}
\usepackage{booktabs}
\usepackage{hyperref}

\title{Replication: The Credit Default Swap–Bond Basis \\
\large Bai, Collin-Dufresne (2019)}
\author{Nicholas Kebo \and Lucie Martin}
\date{\today}

\begin{document}

\maketitle

\section{Introduction}

This document presents a replication of Figure 1 and Table 1 from Bai and Collin-Dufresne (2019), "The Credit Default Swap–Bond Basis." The paper examines the relationship between CDS spreads and bond-implied CDS spreads (PECDS) across different credit rating categories and crisis phases.

The CDS-bond basis is defined as the difference between the market CDS spread and the bond-implied CDS spread:
\begin{equation}
    \text{Basis} = \text{CDS Spread} - \text{PECDS}
\end{equation}

\section{Background and Economic Context}

\subsection{The CDS-Bond Basis Arbitrage}

\subsubsection{The Theoretical Arbitrage}
In a frictionless market, the CDS spread and bond-implied credit spread should be equal. When the basis is negative (CDS $<$ PECDS), an arbitrageur could theoretically:
\begin{enumerate}
    \item Purchase the corporate bond
    \item Finance the purchase via the repo market
    \item Buy CDS protection on the bond
    \item Lock in a risk-free profit equal to the basis
\end{enumerate}
This arbitrage should eliminate any persistent negative basis.

\subsubsection{Why the Arbitrage Failed: Limits to Arbitrage}
However, during the 2007-2009 financial crisis, the basis became persistently and severely negative, suggesting the arbitrage was not risk-free. The paper investigates four key frictions:
\begin{itemize}
    \item \textbf{Collateral Quality}: Bonds with lower credit ratings face higher haircuts in repo markets, requiring more risk capital
    \item \textbf{Liquidity Risk}: Illiquid bonds have higher transaction costs and risk of fire-sale losses
    \item \textbf{Counterparty Risk}: CDS protection is worthless if the protection seller defaults alongside the reference entity
    \item \textbf{Funding Liquidity Risk}: Arbitrageurs face mark-to-market losses when the basis widens as funding costs spike
\end{itemize}

\section{Data Sources}

The replication uses the following data sources:

\begin{itemize}
    \item \textbf{CDS Data}: Markit CDS spreads from WRDS, including 5-year USD-denominated contracts
    \item \textbf{Bond Data}: WRDS Bond Returns database with monthly bond prices and yields
    \item \textbf{Bond Characteristics}: Mergent FISD (Fixed Income Securities Database) for bond issue details
    \item \textbf{Credit Ratings}: Mergent FISD ratings from S\&P and Moody's
\end{itemize}

\textbf{Sample Period}: July 1, 2006 to December 30, 2014

\section{Methodology}

\subsection{Data Preparation}

Bonds are filtered according to the following criteria from the paper:
\begin{enumerate}
    \item Remove convertible bonds
    \item Keep only fixed or zero coupon bonds (no floating rate)
    \item Time-to-maturity between 3 and 7.5 years (to match 5-year CDS liquidity)
    \item Remove bonds with missing prices, yields, or company symbols
    \item Keep only rated bonds (investment grade and high yield)
    \item Keep only corporate bond types (debentures, MTNs, zero coupon)
\end{enumerate}

CDS contracts are filtered to:
\begin{itemize}
    \item USD currency only
    \item 5-year tenor (most liquid maturity)
    \item Sample period: July 2006 - December 2014
\end{itemize}

\subsection{Descriptive Summary of the Underlying Data}

Below are some summary statistics for the underlying matched bond-CDS dataset and the key variables used in the analysis. Table~\ref{tab:sample_summary} reports summary statistics for the main spread variables, while Figure~\ref{fig:underlying_spreads} plots the average daily CDS spread and average daily bond-implied CDS spread proxy (PECDS) over time. Together, these descriptive statistics provide context for the replicated basis results by showing the typical levels, dispersion, and co-movement of the underlying spread series.

\begin{table}[htbp]
    \centering
    \caption{Summary statistics of the underlying matched sample. This table reports descriptive statistics for the key variables used in the analysis: the market CDS spread, the PECDS proxy constructed from bond yields, the CDS--bond basis in basis points, and bond yields. The main takeaway is that the basis is only one part of a broader spread environment, and understanding the levels and dispersion of the underlying series is helpful for interpreting the replicated results.}
    \label{tab:sample_summary}
    \input{../_output/sample_summary_table.tex}
\end{table}

\begin{figure}[htbp]
    \centering
    \includegraphics[width=0.9\textwidth]{../_output/underlying_spreads.png}
    \caption{Average daily CDS spreads and PECDS over time in the matched sample. This figure helps motivate the basis analysis by showing the behavior of the two underlying spread series whose difference defines the CDS--bond basis. Periods in which the two series diverge are precisely the periods that generate large positive or negative basis values.}
    \label{fig:underlying_spreads}
\end{figure}

\subsection{Simplified PECDS Calculation}

This replication uses a simplified approach where the monthly bond yield is used directly as a proxy for PECDS, rather than implementing the full survival probability bootstrapping methodology in Appendix A of the paper. This simplification is necessary due to the lack of available swap rate curve data.

\subsection{Crisis Phases}

Following the paper, the sample is divided into four phases:
\begin{itemize}
    \item \textbf{Before Crisis}: July 2006 - June 2007
    \item \textbf{Crisis I}: July 2007 - August 2008 (subprime crisis to Lehman bankruptcy)
    \item \textbf{Crisis II}: September 2008 - September 2009 (post-Lehman failure)
    \item \textbf{Post-crisis}: October 2009 - December 2014
\end{itemize}

\section{Results}

\subsection{Figure 1: Time Series of CDS-Bond Basis}

Figures~\ref{fig:figure1_replication} and~\ref{fig:figure1_extension} show the dispersion of the CDS--bond basis over time, separated by credit rating. Panel (a) displays investment-grade bonds, while panel (b) shows high-yield bonds. The solid line represents the median basis, while the dashed lines show the 10th and 90th percentiles.

\begin{figure}[htbp]
    \centering
    \includegraphics[width=0.95\textwidth]{../_output/replication_figure1.png}
    \caption{Dispersion of the CDS--bond basis by credit rating (replication window). Panel (a) shows investment-grade bonds; panel (b) shows high-yield bonds. Solid line: median. Dashed lines: 10th and 90th percentiles.}
    \label{fig:figure1_replication}
\end{figure}

\begin{figure}[htbp]
    \centering
    \includegraphics[width=0.95\textwidth]{../_output/extension_figure1.png}
    \caption{Dispersion of the CDS--bond basis by credit rating (extended window). Panel (a) shows investment-grade bonds; panel (b) shows high-yield bonds. Solid line: median. Dashed lines: 10th and 90th percentiles.}
    \label{fig:figure1_extension}
\end{figure}

Key observations from Figures~\ref{fig:figure1_replication} and~\ref{fig:figure1_extension}:
\begin{itemize}
    \item The basis was relatively stable and close to zero before the financial crisis
    \item Both investment-grade and high-yield bonds show sharp negative bases during Crisis II (post-Lehman)
    \item High-yield bonds exhibit greater dispersion and more negative bases than investment-grade bonds
    \item The basis gradually recovered in the post-crisis period but remained negative on average
\end{itemize}

\subsection{Table 1: Summary Statistics}

Table~\ref{tab:table1_replication} presents summary statistics of the CDS-bond basis for the replication period (July 2006 - December 2014), matching the paper's sample. Table~\ref{tab:table1_extension} shows results for the extended period (January 2015 - Present). Statistics include the mean, standard deviation (SD), 10th percentile (P10), and 90th percentile (P90).

\input{../_output/table1_replication.tex}

\input{../_output/table1_extension.tex}

\section{Challenges and Discussion}

\subsection{Successes}
\begin{itemize}
    \item Successfully matched bonds to CDS entities using ticker symbols
    \item Replicated the overall patterns observed in the original paper
    \item Automated the entire data pipeline from raw data to publication-ready outputs
\end{itemize}

\subsection{Challenges and Deviations}
\begin{itemize}
    \item \textbf{Missing MR Clause}: The CDS data contains only XR14 and XR restructuring clauses, not the Modified Restructuring (MR) clause specified in the paper. This is likely due to market convention changes after the sample period.
    
    \item \textbf{Simplified PECDS Calculation}: Due to unavailable swap rate curves, we use bond yields directly as PECDS proxies rather than implementing the full bootstrapping methodology from Appendix A. This may introduce bias as bond yields include both credit spread and risk-free rate components.
    
    \item \textbf{Memory and Performance Issues}: Initial attempts to match bonds with CDS data created a 92-million row dataset from the entire 50-million row CDS universe, causing memory crashes. We resolved this by filtering CDS contracts to 5-year tenor and USD currency before matching rather than after, reducing the matched dataset to approximately 17 million rows.
    
    \item \textbf{Time-to-Maturity in Extended Period}: In the extension period (2015-present), bonds naturally age out of the 3-7.5 year maturity window specified in the paper's methodology. It is unclear whether the extension analysis should use a resampled bond universe for each period or maintain the same bond universe as bonds age, potentially resulting in some bonds with very low time-to-maturity in later years.
    
    \item \textbf{Company Matching}: The paper notes that CDS-bond matching was done "manually," suggesting this is a known difficulty. We use automated ticker matching which may miss some valid pairs or include incorrect matches.
    
    \item \textbf{Sample Coverage}: Our matched sample includes 698 unique CDS names compared to 679 in the paper, suggesting reasonable but not perfect replication of the sample.
\end{itemize}

\subsection{Comparison with Original Results}

Comparing our mean basis values with those reported in the paper reveals both similarities and differences. The overall patterns are consistent: negative bases during crisis periods, with Crisis II showing the most negative values. However, the magnitudes differ, likely due to the simplified PECDS calculation and data differences noted above.

A notable feature in the extended period is the sharp drop in the investment-grade basis visible in Figure~\ref{fig:figure1_extension}, particularly around 2020-2022. This likely reflects the COVID-19 pandemic's impact on credit markets, where unprecedented liquidity interventions and market dislocations may have temporarily widened the basis. Alternatively, it could result from bonds aging out of the 3-7.5 year maturity window in the extended sample, changing the composition of the investment-grade universe over time.

\section{Conclusion}

This replication successfully reproduces the main patterns documented in Bai and Collin-Dufresne (2019), demonstrating negative CDS-bond bases during crisis periods with greater magnitudes for high-yield bonds. While the simplified methodology and data limitations prevent exact numerical replication, the economic insights remain intact: arbitrage relationships between CDS and bond markets can break down significantly during periods of financial stress.

\end{document}